\documentclass[11pt]{article}
\usepackage[margin=1in]{geometry}
\usepackage{amsmath,amssymb,amsthm}
\usepackage{mathtools}
\usepackage{booktabs}
\usepackage[colorlinks=true,linkcolor=blue,citecolor=blue]{hyperref}
\usepackage[capitalise,nameinlink]{cleveref}

\newtheorem{theorem}{Theorem}
\newtheorem{proposition}[theorem]{Proposition}
\newtheorem{lemma}[theorem]{Lemma}
\newtheorem{corollary}[theorem]{Corollary}
\newtheorem{definition}[theorem]{Definition}
\newtheorem{remark}[theorem]{Remark}

\newcommand{\Lean}{Lean}
\newcommand{\code}[1]{\texttt{#1}}
\newcommand{\R}{\mathbb{R}}

\emergencystretch=1.5em

\title{What a layer's output can tell you about its attention:\\
an identifiability threshold, machine-checked}
\author{Fan Li}
\date{}

\begin{document}
\maketitle

\begin{abstract}
Attention heads are usually studied by asking what they compute. We ask the
inverse question: what can be recovered from what they output?

For a single softmax attention head, the layer output determines the attention
weights if and only if the value vectors are affinely independent. Since $n$
affinely independent points require $n \le d+1$, any context longer than the
model width plus one admits more than one attention map producing the same
layer output. We quantify the failure: the weights consistent with a given
output contain an affine subspace of dimension at least $n-(d+1)$. For a
2048-token context in a 768-wide model that is a family of dimension at least
$1279$. A heat map is therefore not ``the'' explanation of a layer's output; it
is one point in a family whose dimension has an explicit lower bound.

The argument rests on structural facts that we also verify: softmax is a
probability distribution, attention outputs are convex combinations of the
values, weights depend only on score differences, and a strictly positive
minimum weight contracts the token spread by an explicit factor --- the
mechanism behind rank collapse.

Everything is formalized in \Lean{}~4 with Mathlib: twenty-six declarations,
no \code{sorry}, axiom set contained in $\{\code{propext},
\code{Classical.choice}, \code{Quot.sound}\}$. The mathematics is elementary
and we say so; the contribution is the question, the threshold, and the
machine-checked record.

\end{abstract}

\tableofcontents

\section{Introduction}

\subsection{Two kinds of claim about a model}

A statement about a trained transformer is one of two things. It is either a
statement about \emph{the architecture} --- true for every admissible choice of
weights --- or a statement about \emph{this model}, established by probing,
ablating or measuring it. The two have different evidence, and they are easy to
confuse in prose.

Most of interpretability studies the second kind: a particular head detects a
particular pattern, a particular direction carries a particular feature. Those
claims are supported by experiments on one network. The first kind is rarer,
because the architecture imposes fewer constraints than one would like: a
softmax head can in principle implement almost any convex combination.

This paper asks what the first kind can buy. Concretely, it studies one
question about a single attention head, with no reference to weights at all.

\subsection{The question}

An attention head reads a score vector $z$ and a family of value vectors
$v_1,\dots,v_n$, and emits
\[
  \mathrm{out}_d \;=\; \sum_{i=1}^{n} p_i \, v_{i,d},
  \qquad
  p_i \;=\; \frac{e^{z_i}}{\sum_j e^{z_j}} .
\]
Everyone studies the map $(z,v)\mapsto \mathrm{out}$. We study the inverse
problem: \emph{given the output, and given the values, is $p$ determined?}

This matters for a familiar practice. Attention maps are drawn as
explanations: ``the model looked here''. If the map is recoverable from what
the layer emits, that practice has a mathematical basis. If it is not, then the
drawn map is one of many maps consistent with exactly the same numbers, and no
amount of care in drawing it will make it canonical.

\subsection{Results}

The answer is a threshold.

\begin{theorem}[Identifiability, \cref{thm:inj,thm:conv}]
Fix values $v : I \to \R^D$. The map $p \mapsto \mathrm{out}$ is injective on
probability vectors if and only if the values are affinely independent.
\end{theorem}

\begin{theorem}[Obstruction, \cref{thm:dim}]
Affine independence of $n$ points in $\R^d$ forces $n \le d+1$. When
$n > d+1$ the values are necessarily dependent, and the weights consistent with
a given output contain an affine subspace of dimension at least $n-(d+1)$
(\cref{thm:kerdim,cor:fiber}).
\end{theorem}

Affine independence is \emph{possible} only for short contexts, $n \le d+1$;
longer contexts are never identifiable, and there the failure has a computable
size. Whether a given short context is identifiable is a property of its
values, not of the architecture.

The remaining results are the scaffolding: the probability and convexity
facts the argument uses, the score-difference invariance and its converse, a
gap-controlled approximation, a contraction statement, and two algebraic facts
about depth. \Cref{sec:structure,sec:dynamic} state the ones we use;
\cref{app:ledger} tabulates all of them, including those not restated here.

\subsection{What is not claimed}

\begin{itemize}
\item \textbf{Not about learned weights.} The scores $z$ and values $v$ are
  free variables, not outputs of projections. The only matrices in the paper
  are the fixed constants of \cref{prop:depth}; they are given, not trained.
  Statements about a named model are outside what is formalized; an earlier
  draft of the specification claimed one and the audit caught it.
\item \textbf{Not new mathematics, mostly.} The contraction is rank collapse
  \cite{dong}, the approximation is the softmax-to-argmax estimate
  \cite{hardattn}, the softmax properties are textbook. The identifiability
  statement we have not found in the literature, but it is elementary linear
  algebra and may well be a remark somewhere; see \cref{sec:related}.
\item \textbf{Not a full transformer.} No multi-head projections, no positional
  encoding, no normalization in the residual stream, no MLP, no causal mask.
\end{itemize}

\section{Setting}\label{sec:setup}

Throughout, $I$ and $D$ are finite nonempty sets; $I$ indexes positions and
$D$ indexes feature coordinates. Nonemptiness is needed because the softmax
denominator is a sum over $I$. A vector is a function $I \to \R$ or $D \to \R$, and
$v : I \to D \to \R$ is a family of $|I|$ vectors in $\R^{|D|}$.

\begin{definition}[Softmax and attention]
For $z : I \to \R$ and $i \in I$,
\[
  \mathrm{softmax}(z)_i = \frac{e^{z_i}}{\sum_{j} e^{z_j}},
  \qquad
  \mathrm{attn}(p,v)_d = \sum_i p_i\, v_{i,d}.
\]
\end{definition}

We write $1$ for the constant function, so that $z + c$ denotes adding the
constant $c$ to every coordinate.

\section{The identifiability threshold}\label{sec:ident}

\subsection{Invertibility}

\begin{definition}[Affine independence, in the form we use]
\label{def:affind}
The values $v : I \to D \to \R$ are affinely independent when the only
coefficient function $c : I \to \R$ with $\sum_i c_i v_{i,d} = 0$ for every $d$
and $\sum_i c_i = 0$ is $c = 0$.
\end{definition}

\begin{theorem}[Identifiability]\label{thm:inj}
If the values are affinely independent and $p, p'$ are probability vectors
($\sum p = \sum p' = 1$) with $\mathrm{attn}(p,v) = \mathrm{attn}(p',v)$, then
$p = p'$.
\end{theorem}

\begin{proof}
The difference $c = p - p'$ satisfies $\sum_i c_i v_{i,d} = 0$ for all $d$, by
subtracting the two outputs, and $\sum_i c_i = 0$ by subtracting the two normalizations.
Affine independence gives $c = 0$.
\end{proof}

\begin{theorem}[Converse]\label{thm:conv}
If the values are affinely dependent, then there exist two distinct probability
vectors $p \ne p'$ with $\mathrm{attn}(p,v) = \mathrm{attn}(p',v)$.
\end{theorem}

\begin{proof}
Take a nontrivial $c$ with $\sum_i c_i v_{i,d} = 0$ for all $d$ and
$\sum_i c_i = 0$, and let $p$ be the uniform weight $1/n$. For $\epsilon > 0$
small enough that $p + \epsilon c \ge 0$ coordinatewise, $p' = p + \epsilon c$
is again a probability vector; its total weight is unchanged because
$\sum c_i = 0$, and its output is unchanged because $c$ annihilates the values.
\end{proof}

\Cref{thm:inj,thm:conv} together give the ``if and only if'' of the
introduction.

\subsection{What the scores contribute}

\begin{proposition}[Converse of translation invariance]\label{prop:shift}
Two softmax distributions coincide, $\mathrm{softmax}(z) = \mathrm{softmax}(z')$,
if and only if $z' = z + c$ for some constant $c$.
\end{proposition}

The forward direction is the familiar shift invariance: adding a constant to
all scores does not change the weights. The content here is the converse, which
says softmax quotients the scores by \emph{exactly} one degree of freedom and no
more. It follows by taking logarithms: $\log \mathrm{softmax}(z)_i = z_i - \log
\sum_j e^{z_j}$, so equality of two softmax vectors forces the two scores to
differ by the difference of the two log-partition functions.

Combining \cref{thm:inj,prop:shift}: for affinely independent values, the layer
output determines the attention weights, and the weights determine the scores up
to the shift that was never observable to begin with.

\subsection{The obstruction}\label{sec:obstruction}

\begin{theorem}[Dimension]\label{thm:dim}
Affine independence of $v : I \to D \to \R$ forces $|I| \le |D| + 1$.
\end{theorem}

\begin{proof}
Consider the linear map
\[
  \Phi : (I \to \R) \to (D \to \R) \times \R, \qquad
  \Phi(c) = \Big(d \mapsto \sum_i c_i v_{i,d},\ \sum_i c_i\Big).
\]
Affine independence says exactly that $\ker \Phi = 0$, i.e.\ that $\Phi$ is
injective. An injective linear map has domain dimension at most codomain
dimension, so $|I| \le |D| + 1$.
\end{proof}

The contrapositive is the statement of interest: once the context is longer
than the width plus one, \emph{no} values are affinely independent, so
\cref{thm:inj} never applies. Long contexts are exactly where attention maps
are most often drawn.

\begin{theorem}[Size of the ambiguity]\label{thm:kerdim}
For any values, $\dim \ker \Phi \ge |I| - (|D| + 1)$.
\end{theorem}

\begin{proof}
Rank--nullity: $\dim(I\to\R) = \dim \operatorname{im}\Phi + \dim \ker \Phi$,
with $\dim (I\to\R) = |I|$ and $\dim \operatorname{im}\Phi \le \dim((D\to\R)\times\R) = |D|+1$.
\end{proof}

\begin{corollary}\label{cor:fiber}
If $c \in \ker \Phi$ and $p$ is any weight vector, then $p + c$ gives the same
output and the same total weight as $p$.
\end{corollary}

So the set of explanations consistent with one output is not a pair of points
but an affine space: of dimension $|I| - \operatorname{rank}\Phi$, and at least
$|I| - |D| - 1$. A 2048-token context in a 768-wide model has a space of
explanations of dimension at least $2048 - 768 - 1 = 1279$.

\begin{remark}
The bound is on the affine space through the explanations, not on its
intersection with the probability simplex. That intersection is a polytope whose
dimension can be smaller --- on the boundary of the simplex it can collapse ---
and computing it is a linear-programming question we do not address.
\end{remark}

\section{Structural facts}\label{sec:structure}

The proof of \cref{thm:inj} uses only that the weights sum to one. The
following are the facts about softmax itself that we verify alongside.

\begin{proposition}[Probability]\label{prop:prob}
$\mathrm{softmax}(z)_i > 0$ for all $i$, and $\sum_i \mathrm{softmax}(z)_i = 1$.
\end{proposition}

\begin{proposition}[Convexity]\label{prop:conv}
If $|v_{i,d}| \le M$ for all $i$, then
$|\mathrm{attn}(\mathrm{softmax}(z),v)_d| \le M$.
\end{proposition}

The proof uses only that softmax weights are nonnegative and sum to one, so the
statement holds for any such $p$; we state it for softmax because that is the
form verified.

\begin{proposition}[Score differences]\label{prop:diff}
$\mathrm{softmax}(z + c) = \mathrm{softmax}(z)$, and $z_i \le z_j$ implies
$\mathrm{softmax}(z)_i \le \mathrm{softmax}(z)_j$.
\end{proposition}

\begin{proposition}[Gap control]\label{prop:gap}
Let $M$ satisfy $|v_{i,d}| \le M$ for all $i$ and $d$, and let $\delta > 0$. If $i_0$ is the unique
maximizer of $z$ and every other score is at least $\delta$ below it, then
\[
  \mathrm{softmax}(z)_{i_0} \ge \frac{1}{1 + (|I|-1)e^{-\delta}} .
\]
Moreover, for \emph{every} $i_0$,
$\bigl|\mathrm{attn}(\mathrm{softmax}(z),v)_d - v_{i_0,d}\bigr| \le
(1 - \mathrm{softmax}(z)_{i_0}) \cdot 2M$.
\end{proposition}

The second half holds without any gap hypothesis; it is the first half that
makes it informative, by forcing $1 - p_{i_0}$ small. The constant $2M$ is
sharp: with two positions, $z = 0$, $v = (2,-2)$ and $M = 2$ the bound is
attained, and replacing $2M$ by $M$ makes it false.

\section{Two dynamical and algebraic facts}\label{sec:dynamic}

\begin{theorem}[Contraction]\label{thm:contract}
Suppose the weight matrix has rows summing to one with every entry at least
$m \ge 0$, and the values are real numbers lying in an interval of length $D$.
Then for any two positions,
\[
  \bigl|\mathrm{attn}(p_i, v) - \mathrm{attn}(p_j, v)\bigr|
  \;\le\; (1 - |I|\,m)\, D ,
\]
\end{theorem}

Softmax weights are strictly positive, so $m > 0$ and the factor is strictly
below one: every layer contracts the spread of the token representations. At
uniform weights $m = 1/|I|$ the factor is zero, i.e.\ the outputs are already
identical; at concentrated weights $m \to 0$ the factor tends to one, i.e.\
hard attention does not contract. The bound interpolates correctly at both ends,
which is a check on the constant.

This is the one-step mechanism behind rank collapse \cite{dong}: absent
anything else, repeated application drives the tokens together. We do not
formalize the iteration; the statement here is for a single layer.

\begin{proposition}[Depth and residuals]\label{prop:depth}
A linear attention layer $X \mapsto P X W$ composes as
$(P_2 W_2)\circ(P_1 W_1) = (P_2 P_1, W_1 W_2)$: with fixed matrices, any depth
collapses to a single layer. With residuals,
$(1+A)(1+B) = 1 + A + B + AB$: the product introduces the cross term $AB$, which
neither factor contains on its own. We do not claim that this makes depth
necessary; see \cref{sec:limits}.
\end{proposition}

The first half is matrix associativity. Its scope is exactly fixed $P$: in real
linear attention the weights depend on the input, and the collapse does not
occur.

\section{Formalization}\label{sec:formal}

The development is in \Lean{}~4 with Mathlib. Every declaration named in the
ledger has three independent checks: the library builds; a script confirms no
\code{sorry}, \code{admit}, \code{native\_decide}, \code{implemented\_by} or
\code{extern} occurs in the sources; and an axiom audit prints the transitive
axiom set of each public declaration and compares it to $\{\code{propext},
\code{Classical.choice}, \code{Quot.sound}\}$.

Two further checks are worth describing because they found things.

\paragraph{Statements are compared, not assumed.} Each result has a challenge
declaration in a separate file, written against the specification; a script
compares the challenge's type with the theorem's type by definitional equality.
The check is deliberately stricter than mathematics --- it rejects a difference
in binder structure that is not a difference in meaning --- and it caught such a
mismatch on its first run.

\paragraph{Independent audit.} Three reviewers with different lenses (statement
fidelity, proof correctness, mathematical content) examined the development,
writing their own test files rather than reading. The correctness reviewer found
no false theorem; the fidelity reviewer found seven overclaims, all in prose
rather than in proofs, including a specification that claimed to formalize a
named frontier architecture when the formalization contains no learned
projection matrices at all. All seven are fixed and recorded.

\section{Related work}\label{sec:related}

The piecewise-linear structure of ReLU networks and the counting of linear
regions are studied in \cite{montufar,telgarsky,arora}. Rank collapse --- the
tendency of pure attention to lose rank with depth --- is established in
\cite{dong}, with a 2024 literature on its spectral mechanism \cite{gap} and on
architectural fixes \cite{lambda}. The sharp Lipschitz constant of softmax is
treated in \cite{lip}. Approximation of hard attention by soft attention is
studied in \cite{hardattn}. The characterization of choice probabilities that leads to the exponential
form of the weights goes back to Luce \cite{luce} and McFadden \cite{mcfadden}.
Two bibliographic qualifications are recorded rather than hidden: the DOI for
Luce refers to a later reissue, not the 1959 edition, and McFadden's text is a
book chapter with no chapter-level DOI.

Explainability methods --- LIME, SHAP, Grad-CAM, integrated gradients
\cite{lime,shap,gradcam,ig} --- produce explanations of a trained model by other
means. This paper is not a method; it asks what information a layer's output
contains about its attention at all.

Searches for attention identifiability and invertibility returned no statement
of \cref{thm:inj}. The closest work \cite{unmask} asks the opposite question:
recovering the input data from the attention weights and the output. Directions
differ; whether the question we ask has appeared as a remark elsewhere we cannot
rule out.

\section{Limitations}\label{sec:limits}

\begin{enumerate}
\item \textbf{A single head; no learned projections.} There are no projections,
  no multiple heads, no positional encoding, no normalization, no MLP. The results are about softmax
  and weighted averages, not about a transformer layer as implemented.
\item \textbf{Elementary mathematics.} The identifiability theorem is two lines;
  the dimension bound is rank--nullity. The contribution is the question and the
  checked record, not depth.
\item \textbf{The simplex intersection is not computed.} \cref{thm:kerdim}
  bounds the affine space containing the explanations; the dimension of its
  intersection with the probability simplex is a linear-programming question we
  leave open.
\item \textbf{Not the trained-model question.} Whether a particular model's
  values are affinely independent is an empirical question about that model. We
  give the criterion; we do not evaluate it on any network.
\end{enumerate}

\appendix

\section{The ledger}\label{app:ledger}

\begin{center}
\begin{tabular}{@{}lll@{}}
\toprule
ID & Content & Role \\
\midrule
C1--C3 & softmax is a probability distribution & support \\
C4 & attention output lies in the range of the values & support \\
C5--C6 & shift invariance, monotonicity & support \\
C7 & gap-controlled top-1 weight & quantitative \\
C8 & distance to any single value $\le (1-p)\,2M$ & quantitative \\
C9 & \code{rms} is homogeneous of degree one & isolated \\
C10--C11 & residual block range and proxy distance & structural \\
C12--C13 & softmax is monotone; weighted variance $\ge 0$ & structural \\
C14--C15 & value path is nonexpansive; block bound & structural \\
C16--C17 & shift invariance forces a Cauchy equation & partial uniqueness \\
C18 & depth hierarchy for attention & \textbf{open, not claimed} \\
C19 & Cauchy classification & \textbf{not formalized} \\
C20--C21 & fixed-$P$ collapse; residual cross term & algebraic \\
C22 & attention contracts the token spread & dynamical \\
C23 & affine independence $\Rightarrow$ identifiability & \textbf{main} \\
C24 & softmax determines scores up to a constant & main \\
C25 & affine independence forces $n \le d+1$ & \textbf{main} \\
C26--C27 & the ambiguity has dimension $\ge n-d-1$ & \textbf{main} \\
C28 & affine dependence yields a second explanation & main \\
\bottomrule
\end{tabular}
\end{center}

\begin{thebibliography}{99}

\bibitem{luce}
R.~D. Luce.
\newblock \emph{Individual Choice Behavior: A Theoretical Analysis}.
\newblock Wiley, 1959. Reissued by Dover, 2005, ISBN 978-0-486-44136-8;
  doi:10.1037/14396-000 refers to the reissue.

\bibitem{mcfadden}
D.~McFadden.
\newblock Conditional logit analysis of qualitative choice behavior.
\newblock In P.~Zarembka, editor, \emph{Frontiers in Econometrics},
  Academic Press, New York, 1974. No chapter-level DOI; the volume's details
  are confirmed by Crossref record doi:10.1017/S0770451800103604.

\bibitem{dong}
Y.~Dong, J.-B. Cordonnier, and A.~Loukas.
\newblock Attention is not all you need: pure attention loses rank doubly
  exponentially with depth.
\newblock arXiv:2103.03404, 2021.

\bibitem{gap}
\newblock Mind the gap: a spectral analysis of rank collapse and signal
  propagation in attention layers.
\newblock arXiv:2410.07799, 2024.

\bibitem{lambda}
\newblock Lambda-skip connections: the architectural component that prevents
  rank collapse.
\newblock arXiv:2410.10609, 2024.

\bibitem{lip}
\newblock Softmax is $1/2$-Lipschitz: a tight bound across all $\ell_p$ norms.
\newblock arXiv:2510.23012, 2025.

\bibitem{hardattn}
\newblock Simulating hard attention using soft attention.
\newblock arXiv:2412.09925, 2024.

\bibitem{unmask}
Y.~Deng, Z.~Song, S.~Xie, and C.~Yang.
\newblock Unmasking transformers: a theoretical approach to data recovery via
  attention weights.
\newblock arXiv:2310.12462, 2023.

\bibitem{montufar}
G.~Mont\'ufar, R.~Pascanu, K.~Cho, and Y.~Bengio.
\newblock On the number of linear regions of deep neural networks.
\newblock \emph{Advances in Neural Information Processing Systems}, 2014.
\newblock arXiv:1402.1869.

\bibitem{telgarsky}
M.~Telgarsky.
\newblock Benefits of depth in neural networks.
\newblock \emph{Conference on Learning Theory}, 2016.
\newblock arXiv:1602.04485.

\bibitem{arora}
R.~Arora, A.~Basu, P.~Mianjy, and A.~Mukherjee.
\newblock Understanding deep neural networks with rectified linear units.
\newblock \emph{International Conference on Learning Representations}, 2018.
\newblock arXiv:1611.01491.

\bibitem{lime}
M.~T. Ribeiro, S.~Singh, and C.~Guestrin.
\newblock ``Why should I trust you?'': explaining the predictions of any
  classifier.
\newblock \emph{ACM SIGKDD}, 2016.
\newblock arXiv:1602.04938.

\bibitem{shap}
S.~M. Lundberg and S.-I. Lee.
\newblock A unified approach to interpreting model predictions.
\newblock \emph{Advances in Neural Information Processing Systems}, 2017.
\newblock arXiv:1705.07874.

\bibitem{gradcam}
R.~R. Selvaraju, M.~Cogswell, A.~Das, R.~Vedantam, D.~Parikh, and D.~Batra.
\newblock Grad-CAM: visual explanations from deep networks via gradient-based
  localization.
\newblock \emph{International Conference on Computer Vision}, 2017.
\newblock arXiv:1610.02391. doi:10.1109/ICCV.2017.74.

\bibitem{ig}
M.~Sundararajan, A.~Taly, and Q.~Yan.
\newblock Axiomatic attribution for deep networks.
\newblock \emph{International Conference on Machine Learning}, 2017.
\newblock arXiv:1703.01365.

\end{thebibliography}

\end{document}
